%% 국문 요약문 (지침 1.다.⑨, 참조 9)
%% 영문논문 의무 제출. 약 800~1,000자 권장.

본 학위논문은 스파이킹 신경망(Spiking Neural Network, SNN)의 직접 학습 과정에서 발생하는 두 가지 핵심 난제, 즉 시간적 공변량 변화(Temporal Covariate Shift, TCS)와 임계전압 학습 시의 기울기 불안정성을 해결하기 위한 새로운 방법을 제안한다.

첫 번째 기여인 \textbf{MP-Init}(Membrane Potential Initialization, 막전위 초기화)은 TCS의 근본 원인이 활성화(activation)가 아닌 LIF 뉴런의 막전위(membrane potential)에 있음을 이론적$\cdot$실증적으로 규명하고, 이를 해결하기 위한 효율적인 초기화 기법이다. 본 논문은 이산 시간 LIF 모델에서 막전위가 마르코프 연쇄(Markov chain)를 형성하며 두블린(Doeblin) 미노라이제이션 조건 아래 유일한 정상분포로 기하적으로 수렴함을 증명하였다. 이 정리에 기반하여, 학습 동안 마지막 시점 $U[T]$의 평균을 층(layer)별 누적 평균으로 추정하여 매 배치 시작 시 막전위를 해당 평균값으로 초기화한다. MP-Init은 추가 파라미터가 거의 없으며, 추론 단계에서 표준 LIF 동작을 그대로 보존하므로 기존 뉴로모픽 하드웨어와 완전히 호환된다.

두 번째 기여인 \textbf{TrSG}(Threshold-Robust Surrogate Gradient, 임계값 강건성 기반 서러게이트 그래디언트)는 임계전압 $V_{\text{thr}}$이 학습 가능한 파라미터일 때 기존 서러게이트 그래디언트가 보이는 임계값 스케일 의존성을 처음으로 분석하고 이를 해결하는 방법이다. 본 논문은 기존 방법을 절대 스케일(AS-SG)과 상대 스케일(RS-SG) 두 계열로 분류하고, 각각이 \emph{기울기 범람/고갈} 또는 \emph{기울기 폭발/소실} 문제를 야기함을 보였다. TrSG는 상대 스케일 활성 영역을 유지하면서 순방향(forward) 경로에서 출력에 임계전압을 곱해 기울기 크기에서 $1/V_{\text{thr}}$ 인자를 상쇄시킨다. 이로써 기울기 크기는 임계값에 무관해지면서도 활성 영역은 임계값에 따라 자연스럽게 적응한다. 추론 시에는 단순한 가중치 재조정을 통해 이진(binary) 스파이크 출력을 회복한다.

CIFAR-10/100, ImageNet, DVS-CIFAR10 등 정적$\cdot$동적 영상 데이터셋에서의 실험 결과, MP-Init과 TrSG의 결합은 표준 LIF 뉴런과 다양한 백본(ResNet-19/34, SEW-ResNet-34, 7-layer CNN)에서 일관된 최첨단 성능을 달성하였다. 또한 Transformer 기반 SNN(SpikeFormer2, QKFormer 등)과 이벤트 기반 객체 검출(COCO-2017) 과제에도 변형 없이 적용 가능하여, 본 방법이 분류 과제를 넘어 폭넓은 응용에 효과적임을 확인하였다.

본 연구는 SNN의 직접 학습 안정성을 한 단계 끌어올림으로써, 에너지 효율적 인공지능 패러다임으로서의 SNN의 실용성을 한층 높이는 데 기여한다.
